\documentclass[12pt,a4paper,oneside]{article}
\usepackage[english]{babel}
\usepackage{amsmath}
\usepackage{amsfonts}
\usepackage{amssymb}
\usepackage{graphicx}
\usepackage{natbib}
\usepackage[left=1in,right=1in,top=1in,bottom=1in]{geometry}
\author{Pieter Verlinden}
\title{Assignment2}

\begin{document}
\begin{titlepage}
\begin{center}
\includegraphics[scale=0.25]{logo.jpg} \\
%\includegraphics[scale=width=0.15\textwidth]{.logo.jpg}\\[1cm]
\textsc{\large Katholieke Universiteit Leuven}\\[4cm]
\textsc{\LARGE Data Structures and Algorithms}\\[0.5cm]
\newcommand{\HRule}{\rule{\linewidth}{0.5mm}}
% Title
\HRule \\[0.5cm]
{ \huge \bfseries Assignment 2 \\ \large{ 8-Puzzle }}\\[0.5cm]

\HRule \\[1.5cm]

%Author and lecturer
\begin{minipage}{0.45\textwidth}
\begin{flushleft} \large
\emph{Author:}
Pieter \textsc{Verlinden}\\

\emph{Student ID:}
r0369547\\

\emph{Study:}
1Ba Informatica

\end{flushleft}

\end{minipage}
\hfill
\begin{minipage}{0.45\textwidth}
\begin{flushleft} \large
\emph{Lecturer:} 
Dr.~Ir.~Philip \textsc{Dutr\'{e}}
\end{flushleft}
\end{minipage}

\vfill

%bottom
{\large April 5, 2013}

\end{center}
\end{titlepage}
\tableofcontents
\newpage
\section{Introduction}
The purpose of this assignment is to study the importance of data-structures, in particular the "priority-queue", by writing a program that solves the 8-puzzle-problem (see section \ref{sec:8puz}). 
\\

This puzzle-problem can be solved by implementing the A* algorithm\cite{A*Wiki}. This algorithm uses priorityfunctions such as the Manhatten priority and the Hamming priority (see sections \ref{subsec:ManPrior} and \ref{subsec:HamPrior}), to determine closest path to the goalstate.
\\

In section \ref{sec:expData}, the priorityfunctions will be discussed and some experimental data will be given. With this data, conclusions will be made and questions can be answered in section \ref{sec:QnA}.
\section{8-puzzle}\label{sec:8puz}
\section{Experimental data}\label{sec:expData}
\subsection{Manhatten priority}\label{subsec:ManPrior}
\subsection{Hamming priority}\label{subsec:HamPrior}
\section{Q{\&}A: Questions and Answers}\label{sec:QnA}
	\cite{princetonPuzzle}
\section{Conclusion}
\newpage
\bibliographystyle{plain}
\bibliography{ref}
\end{document}